\honest{Politics Is the Mind-Killer}{Eliezer Yudkowsky}{2007}

People think badly about politics. The reason is ``so obvious as to be worth belaboring'': in the ancestral environment, being on the wrong side of an argument could get you killed, and the right side could let you kill a rival. Arguing about the minimum wage today, you run on those adaptations. I give no evidence for this.

My advice is this. When you want to teach a point about rationality, do not take your example from current politics; use Louis XVI. Politics is an important place to apply your own reasoning, but a terrible place to learn rationality, ``unless all the discussants are already rational.'' \nb{Any group that wants to argue about politics can decide it qualifies.}

Politics is war by other means, and arguments are soldiers. Once you know your side, you must support all its arguments and attack all the other side's. Otherwise you are stabbing your own soldiers in the back. Scientists who weigh all sides at work turn into slogan-chanting zombies when a political issue comes up.

My example is a standard problem from artificial intelligence: all Quakers are pacifists, all Republicans are not, and Nixon is both. I state it with ``all'', which turns it into a plain contradiction. \nb{The problem was invented to show how to reason with rules that have exceptions, and the usual versions state both rules as defaults that allow exceptions.}

I spend two of my eight paragraphs asking why the problem's author, whom I do not name, chose Nixon. Was it to make Republicans feel unwelcome in courses on artificial intelligence? My answer is that it was ``probably'' to get in a dig at the other side. A solid punch feels good, like a chocolate cookie, and, as with cookies, not everything pleasant is good for you. \nb{The post offers no evidence for the motive, and does not consider the plain reason: Nixon was a famous man who was both a Quaker and a Republican.} A footnote assures you that I am neither a Republican nor a Democrat.

My rule is modest. I am not asking you to be apolitical, or even to follow Wikipedia's Neutral Point of View, only to skip gratuitous jabs at a party. If your topic is attempts to ban evolution in schools, discuss it, but do not blame the whole Republican Party; some readers may be Republicans who see the problem as a few rogues. Whether the party really is at fault does not matter. What matters is ``the spiritual growth of the community.'' \nb{So this rule exists to keep the peace. Keeping the peace and finding out what is true are different aims, and here the post argues only for the first.}

\nb{The title says that politics kills minds; the essay says that politics is a poor place to practice. The title is the part that gets quoted. Seven years later, in a discussion on LessWrong, Yudkowsky said the post had meant ``national politics as seen on TV.''}
